\section{Fine-Grained Timing Sensors on FPGAs}
\label{sec:bg-tdc}

Recovering FPGA pentimenti requires the measurement of picosecond-scale
shifts in propagation delay across individual routing paths. FPGAs are
an excellent platform for these measurements because their
reconfigurable fabric supports the placement of dedicated
time-to-digital converter (TDC) sensors connected directly to the routes
being probed. A TDC is a circuit that converts the time difference
between two signal edges into a digital count proportional to the delay
between edges two~\cite{henzler2010time}. In the FPGA security
literature, TDCs have been developed extensively as on-chip
voltage-fluctuation sensors for remote power analysis
attacks~\cite{schellenberg2018remote, zhao2018fpga, gnad2017voltage,
	gnad2016analysis}. Pentimenti-based attack classes repurposes the same
sensor architecture to recover a distribution of propagation delay
measurements induced by BTI rather than co-tenant power distribution
network fluctuations ~\cite{drewes2024pentimento}.

\subsection{TDC Architecture}

The dominant TDC architecture in FPGA security work combines a carry
chain or programmable delay line with a sampling register
array~\cite{gnad2016analysis, schellenberg2018remote, zhao2018fpga,
	gnad2017voltage}. A reference clock edge propagates through a chain of
FPGA primitives, typically the dedicated carry-propagation logic, and a
snapshot of the chain's state is captured by a register array clocked by
a phase-offset version of the same reference. The captured pattern of
high and low values encodes the position of the propagating edge at the
sampling instant, producing a quantized measurement of the delay between
the two clock edges. The number of registers that have captured the high
level before the edge arrives is the digital count proportional to the
delay~\cite{gnad2016analysis, henzler2010time}.

Two parameters determine the achievable measurement resolution. The
\emph{per-stage} resolution is set by the propagation delay of the FPGA
primitive composing the delay line: dedicated carry primitives typically
provide tens of picoseconds per stage in modern FPGA process
nodes~\cite{schellenberg2018remote, drewes2024turn}. The \emph{tunable}
resolution is set by the phase granularity at which the sampling clock
can be offset from the propagating clock; this phase offset is supplied
by the FPGA's phase-locked loops (PLLs) and clock-management primitives.
A TDC's effective resolution is the composition of these two parameters.
A single carry-chain measurement quantizes delay at the per-stage
granularity, while repeated measurements with varying phase offsets
compose a sweep of finer effective
resolutions~\cite{drewes2024pentimento, drewes2024turn}.

\subsection{PLL Phase Shift and Vendor Asymmetry}

The phase granularity available to a TDC's sampling clock is determined
by the configuration capabilities of the FPGA's
PLLs~\cite{drewes2024turn}. Xilinx FPGAs expose a fine-grained
phase-shift primitive on their mixed-mode clock managers (MMCMs) that
permits sub-degree adjustments of the output phase, providing on the
order of $1/56$ of the voltage-controlled-oscillator (VCO) period per
phase step~\cite{drewes2024pentimento}. For typical VCO operating
frequencies, this yields picosecond-level effective TDC resolution
sufficient to resolve BTI-induced delay shifts. Intel FPGAs, by
contrast, expose phase shifting only at the granularity of the VCO
period divided by a small fixed integer in the case of the Cyclone~V
family, $T_\text{VCO}/8$~\cite{drewes2024turn}. This coarser default
granularity is insufficient to resolve single-picosecond BTI signatures
at typical VCO frequencies, and has constrained demonstrations of FPGA
pentimenti recovery to Xilinx products.

PLL output frequencies on Intel platforms are generated by
multiplier-divider structure, and the phase of any given output
frequency can be modulated by selecting alternative multiplier and
divider configurations that produce the same nominal output frequency
through different VCO operating points~\cite{drewes2024turn}. By
enumerating valid configurations within the documented operating ranges
of the PLL, a composite phase sweep at granularity finer than the
default $T_\text{VCO}/8$ can be assembled from individual PLL
configurations, each contributing a single phase point at coarse default
granularity. Chapter~\ref{chap:fpga-sensor} of this dissertation
formalizes this approach as a constrained PLL configuration search and
demonstrates its application on Intel Cyclone~V FPGAs.

\subsection{Pentimenti Measurement on FPGAs}

The TDC architecture described above supports two distinct measurement
regimes. In the power side-channel regime, the sensor is sampled
continuously while a victim workload executes, producing power-trace
measurements suitable for correlation power analysis (CPA) and related
attacks~\cite{schellenberg2018remote, zhao2018fpga,
	clavier14_practical}. In the pentimenti regime, the sensor measures the
propagation delay of a configured route while the fabric is otherwise
idle, producing measurements suitable for recovering BTI-induced delay
shifts after a victim workload has completed. An adversary co-located on
a target FPGA after a victim has released the fabric configures a TDC
connected to a previously-stressed routing net and measures its
propagation delay after applying their own known stress. The measured
route either responds with a characteristic stress pattern or a
characteristic recovery pattern. With knowledge of their own applied
stress the adversary can distinguish if their stress reinforces or
recovers the stress from a prior user \cite{drewes2024pentimento}. Due
to the asymmetries bewteen the burn-in (stress) and recovery curves
between logical 1 and logical 0, the adversary is able to determine the
logical voltage level applied to the route.

BTI-induced delay shifts are on the order of single picoseconds, while
power-supply-induced fluctuations typically span tens to hundreds of
picoseconds. The vendor asymmetry in PLL phase-shift capabilities is
therefore most severe for pentimenti recovery specifically, and not
necessarily for the broader class of FPGA-side timing attacks that have
been deployed across both vendors. Chapter~\ref{chap:fpga-sensor} of
this dissertation closes the asymmetry for pentimenti measurement,
broadening the FPGA platforms on which the attack class can be studied.